% Options for packages loaded elsewhere
\PassOptionsToPackage{unicode}{hyperref}
\PassOptionsToPackage{hyphens}{url}
\PassOptionsToPackage{dvipsnames,svgnames,x11names}{xcolor}
%
\documentclass[
  letterpaper,
  DIV=11,
  numbers=noendperiod]{scrartcl}

\usepackage{amsmath,amssymb}
\usepackage{iftex}
\ifPDFTeX
  \usepackage[T1]{fontenc}
  \usepackage[utf8]{inputenc}
  \usepackage{textcomp} % provide euro and other symbols
\else % if luatex or xetex
  \usepackage{unicode-math}
  \defaultfontfeatures{Scale=MatchLowercase}
  \defaultfontfeatures[\rmfamily]{Ligatures=TeX,Scale=1}
\fi
\usepackage{lmodern}
\ifPDFTeX\else  
    % xetex/luatex font selection
\fi
% Use upquote if available, for straight quotes in verbatim environments
\IfFileExists{upquote.sty}{\usepackage{upquote}}{}
\IfFileExists{microtype.sty}{% use microtype if available
  \usepackage[]{microtype}
  \UseMicrotypeSet[protrusion]{basicmath} % disable protrusion for tt fonts
}{}
\makeatletter
\@ifundefined{KOMAClassName}{% if non-KOMA class
  \IfFileExists{parskip.sty}{%
    \usepackage{parskip}
  }{% else
    \setlength{\parindent}{0pt}
    \setlength{\parskip}{6pt plus 2pt minus 1pt}}
}{% if KOMA class
  \KOMAoptions{parskip=half}}
\makeatother
\usepackage{xcolor}
\setlength{\emergencystretch}{3em} % prevent overfull lines
\setcounter{secnumdepth}{-\maxdimen} % remove section numbering
% Make \paragraph and \subparagraph free-standing
\makeatletter
\ifx\paragraph\undefined\else
  \let\oldparagraph\paragraph
  \renewcommand{\paragraph}{
    \@ifstar
      \xxxParagraphStar
      \xxxParagraphNoStar
  }
  \newcommand{\xxxParagraphStar}[1]{\oldparagraph*{#1}\mbox{}}
  \newcommand{\xxxParagraphNoStar}[1]{\oldparagraph{#1}\mbox{}}
\fi
\ifx\subparagraph\undefined\else
  \let\oldsubparagraph\subparagraph
  \renewcommand{\subparagraph}{
    \@ifstar
      \xxxSubParagraphStar
      \xxxSubParagraphNoStar
  }
  \newcommand{\xxxSubParagraphStar}[1]{\oldsubparagraph*{#1}\mbox{}}
  \newcommand{\xxxSubParagraphNoStar}[1]{\oldsubparagraph{#1}\mbox{}}
\fi
\makeatother


\providecommand{\tightlist}{%
  \setlength{\itemsep}{0pt}\setlength{\parskip}{0pt}}\usepackage{longtable,booktabs,array}
\usepackage{calc} % for calculating minipage widths
% Correct order of tables after \paragraph or \subparagraph
\usepackage{etoolbox}
\makeatletter
\patchcmd\longtable{\par}{\if@noskipsec\mbox{}\fi\par}{}{}
\makeatother
% Allow footnotes in longtable head/foot
\IfFileExists{footnotehyper.sty}{\usepackage{footnotehyper}}{\usepackage{footnote}}
\makesavenoteenv{longtable}
\usepackage{graphicx}
\makeatletter
\newsavebox\pandoc@box
\newcommand*\pandocbounded[1]{% scales image to fit in text height/width
  \sbox\pandoc@box{#1}%
  \Gscale@div\@tempa{\textheight}{\dimexpr\ht\pandoc@box+\dp\pandoc@box\relax}%
  \Gscale@div\@tempb{\linewidth}{\wd\pandoc@box}%
  \ifdim\@tempb\p@<\@tempa\p@\let\@tempa\@tempb\fi% select the smaller of both
  \ifdim\@tempa\p@<\p@\scalebox{\@tempa}{\usebox\pandoc@box}%
  \else\usebox{\pandoc@box}%
  \fi%
}
% Set default figure placement to htbp
\def\fps@figure{htbp}
\makeatother

\KOMAoption{captions}{tableheading}
\makeatletter
\@ifpackageloaded{caption}{}{\usepackage{caption}}
\AtBeginDocument{%
\ifdefined\contentsname
  \renewcommand*\contentsname{Table of contents}
\else
  \newcommand\contentsname{Table of contents}
\fi
\ifdefined\listfigurename
  \renewcommand*\listfigurename{List of Figures}
\else
  \newcommand\listfigurename{List of Figures}
\fi
\ifdefined\listtablename
  \renewcommand*\listtablename{List of Tables}
\else
  \newcommand\listtablename{List of Tables}
\fi
\ifdefined\figurename
  \renewcommand*\figurename{Figure}
\else
  \newcommand\figurename{Figure}
\fi
\ifdefined\tablename
  \renewcommand*\tablename{Table}
\else
  \newcommand\tablename{Table}
\fi
}
\@ifpackageloaded{float}{}{\usepackage{float}}
\floatstyle{ruled}
\@ifundefined{c@chapter}{\newfloat{codelisting}{h}{lop}}{\newfloat{codelisting}{h}{lop}[chapter]}
\floatname{codelisting}{Listing}
\newcommand*\listoflistings{\listof{codelisting}{List of Listings}}
\makeatother
\makeatletter
\makeatother
\makeatletter
\@ifpackageloaded{caption}{}{\usepackage{caption}}
\@ifpackageloaded{subcaption}{}{\usepackage{subcaption}}
\makeatother

\usepackage{bookmark}

\IfFileExists{xurl.sty}{\usepackage{xurl}}{} % add URL line breaks if available
\urlstyle{same} % disable monospaced font for URLs
\hypersetup{
  pdftitle={ORCA: Observability via Declarative Dataflows},
  pdfauthor={CMU and LANL},
  colorlinks=true,
  linkcolor={blue},
  filecolor={Maroon},
  citecolor={Blue},
  urlcolor={Blue},
  pdfcreator={LaTeX via pandoc}}


\title{ORCA: Observability via Declarative Dataflows}
\author{CMU and LANL}
\date{2025-05-14}

\begin{document}
\maketitle


\section{ORCA: Overview}\label{orca-overview}

ORCA (\emph{Observability with Realtime Control and Aggregation}) is
designed for continuous, interactive, and realtime visibility into a
\emph{bulk-synchronous parallel (BSP)} application. ORCA targets
enabling the following workflows:

\begin{enumerate}
\def\labelenumi{\arabic{enumi}.}
\tightlist
\item
  Describe observability intent (a view of application state) via a set
  of SQL queries (called a \texttt{flow}). ORCA orchestrates just the
  necessary data movement needed to materialize this view, using SQL
  operator placement.
\item
  Dynamically update this intent, by installing a new \texttt{flow}. All
  subsequent telemetry is aggregated using the new flow.
\end{enumerate}

These two capabilities enable an entire spectrum of observability needs:
from common-case monitoring (minimal collection of summary statistics)
to root-cause localization if and when a problem is suspected.

\section{Motivation: Lessons from Optimizing Parallel
Codes}\label{motivation-lessons-from-optimizing-parallel-codes}

We spent 2+ years understanding parallel code performance at upto 4,096
cores on an in-house cluster and optimizing it. ORCA emerged from this
experience. Key lessons:

\begin{enumerate}
\def\labelenumi{\arabic{enumi}.}
\item
  Global collectives impose extreme performance constraints: a single
  straggler at a global barrier can cause underutilization of the entire
  system (see Fig. 1).
\item
  Sampling and profiling, popular aggregation approaches in existing
  tooling, obscure problem structure (see Figs. 1 and 2).
\item
  Fine-grained telemetry is currently only available via tracing.
  Tracing formats are not amenable to ad-hoc analytics, and their
  post-mortem nature delays problem localization.

  \newpage
\end{enumerate}

The following needs, thus emerged, for an ideal observability system:

\begin{enumerate}
\def\labelenumi{\arabic{enumi}.}
\tightlist
\item
  Solving performance problems requires synthesizing meaningful views
  from raw telemetry. eBPF enables this at a single-node, but not at
  cluster-scale.
\item
  View definition interface should be programmable and dynamically
  updatable. Even with tracing, it is not possible to collect \emph{all}
  telemetry, nor is it necessary.
\item
  Recurrent performance problems require a top-down analysis, to be able
  to go from \emph{``is my app slow?''} to \emph{``what in the stack is
  causing this?''}.
\end{enumerate}

\begin{figure}[H]

{\centering \includegraphics[width=0.5\linewidth,height=\textheight,keepaspectratio]{amr-figs/lesson_volfunc_v2_short_zoom.pdf}

}

\caption{Simplified view of telemetry from a real BSP app with a
volatile function followed by a global collective. Tail latency in
volatile function gets misattributed as elevated collective time in
aggregate telemetry (left), but not in fine-grained telemetry (right).
Localization revealed an issue in the network stack, mitigating which
eliminated the spikes, and therefore the collective time.}

\end{figure}%%
\begin{figure}[H]

{\centering \includegraphics[width=0.5\linewidth,height=\textheight,keepaspectratio]{tikz-pdf-out/sampling.pdf}

}

\caption{A visualization of the two dimensions of BSP telemetry:
timesteps (epochs), and ranks (workers). Sampling collects an arbitrary
subset of telemetry. Coordinated collection is necessary to understand
problem structure.}

\end{figure}%

We found BSP applications to be highly amenable to SQL-based aggregation
(\texttt{GROUPBY} timestep or epoch, compute summary statistics). We
decided to model telemetry as relational data, and leverage modern OLAP
developments: Arrow, Parquet, and Datafusion to build an observability
engine using in-situ in-memory aggregation for dynamic dataflows.

\section{ORCA: Aggregation and
Control}\label{orca-aggregation-and-control}

ORCA leverages an overlay network to mediate between an active BSP
application and an asynchronous controller. The overlay augments the
controller with a tier of dedicated aggregator nodes (collectively
comprising the \emph{asynchronous/monitoring domain}), which monitors a
BSP application (Fig. 3). Thus a tree topology is created, with leaf
nodes being MPI ranks, which are sharded and assigned to aggregators. An
intermediate tier of lead MPI ranks can be used for lightweight
pre-aggregation to manage fanout at extreme scales.

\begin{figure}[H]

{\centering \includegraphics[width=0.8\linewidth,height=\textheight,keepaspectratio]{tikz-pdf-out/overlay-agg-new.pdf}

}

\caption{An overlay example showing the aggregation flow between MPI
ranks and the monitoring tier, going to controller via MPI reps and
aggregators.}

\end{figure}%

The overlay communicates via Mercury, an RPC library for
high-performance fabrics. ORCA is amenable to all its communication
being relegated to an asynchronous low-QoS tier, so as to not interfere
with application performance. Observability capabilities are built on
two key functions:

\begin{enumerate}
\def\labelenumi{\arabic{enumi}.}
\tightlist
\item
  Bottom-up aggregation in a manner defined by the active \texttt{flow},
  with operator pushdown to minimize data movement, using the overlay as
  a reduction tree.
\item
  Control semantics, enabled by an asynchronous, timestep-consistent
  two-phase commit protocol (\texttt{TS2PC}). \texttt{TS2PC} enables
  replicated state machine semantics on the MPI tier, using the overlay
  as a broadcast tree. Brief \texttt{TS2PC} description:

  \begin{itemize}
  \tightlist
  \item
    \emph{Correctness guarantee}: control commands commit on all ranks
    or none.
  \item
    \emph{Performance guarantee}: Protocol execution is asynchronous and
    non-blocking. If the controller thinks that the application is at
    timestep \(T\), a control message is issued to be applied at
    timestep \(T + \Delta\). \(\Delta\) is a function of overlay latency
    and is learned automatically.
  \item
    \emph{Usage}: \texttt{TS2PC} is used to issue commands such as
    \texttt{activate/deactivate\ probe}, \texttt{pause/resume}
    application, \texttt{update\ flow} etc.
  \end{itemize}
\end{enumerate}

\section{Design Overview}\label{design-overview}

\begin{enumerate}
\def\labelenumi{\arabic{enumi}.}
\tightlist
\item
  \emph{Data collectors} represent emittable streams with different
  schemas. We currently have data collectors for MPI collectives and
  Kokkos, and more can be added for \texttt{NCCL}, \texttt{Pytorch},
  eBPF, Infiniband etc. Each data collector may emit different probes,
  each of which can be individually toggled at runtime. If probes are
  disabled or stream is not needed, the collection call becomes a no-op.
\item
  Data collectors are flushed at the end of each timestep to generate
  Arrow data, which is collected and aggregated via the tree.
  Aggregation also sequences collected data by timestep, to ensure that
  a prior timestep is completed before the next timestep is processed.
\item
  Each layer applies assigned partial aggregations, and pushes the
  results up the tree. Finally, data is collected to different
  \emph{sinks}: these could be on-disk Parquet files, timeseries
  databases serving Grafana etc.
\end{enumerate}

ORCA is written in C++, with query planning/execution in Rust using
Apache Datafusion.

\section{Potential Extensions}\label{potential-extensions}

We briefly describe some potential extensions to ORCA. A demonstration
of the first two is in our roadmap.

\begin{enumerate}
\def\labelenumi{\arabic{enumi}.}
\item
  \emph{Online operator-in-the-loop tuning}: the \texttt{TS2PC}
  mechanism can be used to update various application parameters
  (thresholds, rates etc.) at runtime (must be exposed via ORCA), and
  observe their impact in realtime. ORCA can also form the
  infrastructure hosting autotuning algorithms, providing them control
  and feedback.
\item
  \emph{Observing application state}: in addition to performance
  telemetry, ORCA can also aggregate internal application state (such as
  statistics on particle energies or model weights and gradients).
  Approximation techniques such as sketches can be used to define lossy
  aggregations, to reduce data volume with minimal impact on fidelity.
\item
  \emph{Leveraging accelerators}: ORCA can be extended to use
  accelerators as intermediate nodes. Their capabilities can be exposed
  as special nodes in the dataflow, and the planner can be extended to
  leverage them. Datafusion's architecture enables continuous evolution
  of query planning and execution.
\end{enumerate}




\end{document}
